\documentclass[
	ngerman,							%deutsch
	a4paper,							%Seiten Format
	11pt,							%Schriftgröße
	oneside							%einseitig
	]{article}							%Artikel
	

\usepackage[utf8]{inputenc}				%Zeichenkodierung
\usepackage[T1]{fontenc}				%Kein Plan
\usepackage{lmodern}
\usepackage{fontawesome}
\usepackage[colorlinks=true, linkcolor=ColorTwo, urlcolor=ColorTwo, citecolor=ColorOne]{hyperref}
\usepackage{tikz}
\usepackage{hyperref}
\usepackage{apacite}					%Fürs Zitieren - wichtig muss nach hyperref geladen werden!
\usepackage{graphicx}					%Irgendwas mit Bildern?
\usepackage[english]{babel}			%englische Silbentrennung
\usepackage{pdflscape}


\usepackage{float}						%Um Text um Floats rumfloaten zu lassen
\usepackage{
	amsmath,
	amssymb,
	amsfonts,
	fancyhdr}						%Mathezeugs?
\usepackage{setspace}
\usepackage[left=2.6cm, right=2.9cm,top=3.5cm, bottom=3.5cm]{geometry}
\usepackage{rotating}
\usepackage[labelfont=bf]{caption}
%\usepackage{tocbibind}
%\usepackage{arial}
%\usepackage{helvet}
%\renewcommand{\familydefault}{\sfdefault}
\usepackage[section]{placeins}

%\usepackage{natbib}
%\usepackage[usenames,dvipsnames]{xcolor}

\usepackage{tocloft}

\usepackage{enumitem}
\usepackage{sectsty}
\usepackage{wasysym}
\usepackage{multicol}
\usepackage{tabularx}
\usepackage{adjustbox}
\geometry{margin=1in}

\makeatletter
\setlength{\@fptop}{0pt}
\makeatother
	
\pagestyle{fancy}	

\renewcommand{\BAstyle}{\scshape} 				%Macht \cite{} irgendwie in kapitaelchen
\newcommand{\centered}[1]{\begin{tabular}{l} #1 \end{tabular}}
\usepackage{listings}

\fancypagestyle{portraitstyle}{
  \fancyhf{} % Lösche alle Kopf- und Fußzeilen
  \lhead{Kopfzeile links}
  \rhead{Kopfzeile rechts}
  \cfoot{Seitennummerierung}
}

\title{Summary}  								%Titel
\author{Tobias Manzel} 							%Name
\date{\today}										%Datum
\fancyhead[RE,RO]{Tobias Manzel}						%Text Rechts im Header
\fancyhead[LE,LO]{\leftmark}	 	%Text Links im Header (Hier Kapitel)
% !TEX root = Simple-CV.tex
%-------------------------------------------------------------------------------------------------------
% Packages
%-------------------------------------------------------------------------------------------------------



%-------------------------------------------------------------------------------------------------------
% Layout
%-------------------------------------------------------------------------------------------------------
\pagenumbering{gobble}
\renewcommand{\baselinestretch}{1.5} 
\setlength{\parindent}{0pt}

%
% Color theme
%
\definecolor{ColorOne}{RGB}{1,55,130} 	% Blue
\definecolor{ColorTwo}{RGB}{120,0,0} 	% Mauve
%\definecolor{ColorTwo}{RGB}{140,100,0} 	% Gold

\sectionfont{\color{ColorOne}} 
\subsectionfont{\color{ColorOne}} 

%
% Vertical line
%
\newcommand{\MyVerticalRule}{%
	\textcolor{ColorOne}{\rule{1pt}{\textheight}}
}

%
% Update
%
\newcommand{\LastUpdate}{%
\vfill
\centering \small
\textcolor{ColorOne}{Last updated: \monthname,~\the\year.}
}

%
% Skip
%
\newcommand{\MySkip}{
\vskip12pt
}

%
% Format hyperrefs
%
\newcommand{\myhref}[2]{%
\href{#1}{\textcolor{ColorTwo}{#2}}
}
%
% Format skill bullets
%
\newcommand{\SkillBull}[1]{%
\textcolor{ColorTwo}{#1}
}




\begin{document}

\begin{titlepage}
\centering
%\begin{flushleft}
	\makebox[0pt][l]{%
 	\raisebox{-\totalheight}[0pt][0pt]{%
      \includegraphics[width=0.7\textwidth]{signet.png}}}%

	%\end{flushleft}

	{\bfseries\large University of Greifswald \par}
	{\bfseries\large Faculty of Mathematics and Natural Sciences \par}
	{\bfseries\large Institute for Geographie and Geology \par}
	\vspace{0.25cm}
    {\rule{15cm}{0.4pt} \par}
	\vspace{0.5cm}
	{\scshape\large Code Documentation\par}
	\vspace{0.25cm}
	{\Large\bfseries Full Documentation of AlphaBetaSquared Class
          for Python\par}
	\vspace{1cm}
	{\Large\bfseries Greifswald \today \par}
	\vspace{2cm}
    %{\includegraphics[width=1\textwidth]{title.jpg} \par}
    	\vspace{6cm}
\end{titlepage}

\newpage
\setstretch{1}
\thispagestyle{empty}
\tableofcontents
\newpage

\section{Overview}
The \texttt{AlphaBetaSquared} class is a Python utility designed for
statistical analysis and visualization of particle size distributions
from CSV data files. It computes key statistical parameters ($\alpha$
and $\beta^2$), fits log-normal distributions, and generates plots to
visualize both size distributions and kinetic analysis after
\shortciteNP{eberl1998deducing} (alpha-beta squared plots).\\
\ \\
\noindent
The class supports flexible data loading, customizable plotting, and data export functionalities, making it suitable for both interactive use (e.g., Jupyter notebooks) and batch processing.

\section{Class Description}
\subsection*{Class Definition}
\begin{lstlisting}
class AlphaBetaSquared():
\end{lstlisting}

\subsection*{Purpose}
The \texttt{AlphaBetaSquared} class provides methods to:
\begin{itemize}
    \item Load data from one or more CSV files.
    \item Calculate statistical metrics ($\alpha$, $\beta$, maximum values, and log-normal fits).
    \item Visualize particle size distributions and $\alpha$ vs. $\beta^2$ plots.
    \item Export computed data to CSV files.
\end{itemize}

\subsection*{Attributes}
\begin{itemize}
    \item \texttt{self.data}: Dictionary storing loaded CSV data.
    \item \texttt{self.\_alpha}: Dictionary of computed $\alpha$ values.
    \item \texttt{self.\_beta}: Dictionary of computed $\beta$ values.
    \item \texttt{self.\_table\_max}: Dictionary of maximum values per column.
    \item \texttt{self.\_lognorm\_fits}: Dictionary of log-normal fit parameters.
    \item \texttt{self.auto\_display}: Boolean to control automatic plot display (default: \texttt{False}).
    \item \texttt{self.alphabeta\_scale}: Boolean to scale $\alpha$-$\beta^2$ plot limits (default: \texttt{False}).
    \item \texttt{self.save\_plot}: Boolean to save plots to files (default: \texttt{False}).
    \item \texttt{self.plot\_config}: Dictionary for plot formatting (partially implemented).
    \item \texttt{self.\_\_exportable\_attributes}: Tuple of attributes available for export.
\end{itemize}
\newpage
\subsection*{Dependencies}
\begin{table}[ht]
\centering
\begin{tabular}{lll}
\textbf{Package} & \textbf{Version} & \textbf{Installation Command} \\
Python & $\geq$3.8 & \texttt{conda install python=3.9.2} \\
NumPy & 1.23.x & \texttt{pip install numpy==1.23.4} \\
SciPy & 1.9.x & \texttt{pip install scipy==1.9.1} \\
Matplotlib & 3.6.x & \texttt{pip install matplotlib==3.6.2} \\
Pandas & 2.1.x & \texttt{pip install pandas==2.1.4} \\
\end{tabular}
\end{table}

\section{Methods}

\subsection{\texttt{\_\_init\_\_}}
\begin{lstlisting}
def __init__(self, *args, plot_config=None)
\end{lstlisting}
\textbf{Description:} Initializes the class by loading data and computing statistical metrics.

\textbf{Parameters:}
\begin{itemize}
    \item \texttt{*args}: Variable-length argument list of CSV file paths or a single list of paths.
    \item \texttt{plot\_config}: Optional dictionary to customize plot appearance (default: \texttt{None}).
\end{itemize}

\textbf{Behavior:} Loads data, calculates $\alpha$, $\beta$, maximum values, and log-normal fits, and sets up plotting attributes.

\subsection{\texttt{\_load\_data}}
\begin{lstlisting}
def _load_data(self, args)
\end{lstlisting}
\textbf{Description:} Loads data from CSV files into a dictionary.

\textbf{Parameters:}
\begin{itemize}
    \item \texttt{args}: Tuple of file paths or a single list of paths.
\end{itemize}

\textbf{Returns:} Dictionary with keys as table names (derived from filenames) and values as \texttt{pandas.DataFrame} objects.

\subsection{\texttt{set\_auto\_display}}
\begin{lstlisting}
def set_auto_display(self, x)
\end{lstlisting}
\textbf{Description:} Sets the \texttt{auto\_display} attribute.

\textbf{Parameters:}
\begin{itemize}
    \item \texttt{x}: Boolean value.
\end{itemize}

\subsection{\texttt{set\_save\_plot}}
\begin{lstlisting}
def set_save_plot(self, x)
\end{lstlisting}
\textbf{Description:} Sets the \texttt{save\_plot} attribute.

\textbf{Parameters:}
\begin{itemize}
    \item \texttt{x}: Boolean value.
\end{itemize}

\subsection{\texttt{set\_alphabeta\_scale}}
\begin{lstlisting}
def set_alphabeta_scale(self, x)
\end{lstlisting}
\textbf{Description:} Sets the \texttt{alphabeta\_scale} attribute.

\textbf{Parameters:}
\begin{itemize}
    \item \texttt{x}: Boolean value.
\end{itemize}

\subsection{\texttt{calc\_alpha}}
\begin{lstlisting}
def calc_alpha(self, x)
\end{lstlisting}
\textbf{Description:} Calculates the mean of the natural logarithm of the input data.

\textbf{Parameters:}
\begin{itemize}
    \item \texttt{x}: Numeric array-like object.
\end{itemize}

\textbf{Returns:} Float representing $\alpha$.

\subsection{\texttt{calc\_beta}}
\begin{lstlisting}
def calc_beta(self, x)
\end{lstlisting}
\textbf{Description:} Calculates the variance of the natural logarithm of the input data.

\textbf{Parameters:}
\begin{itemize}
    \item \texttt{x}: Numeric array-like object.
\end{itemize}

\textbf{Returns:} Float representing $\beta$.

\subsection{\texttt{show}}
\begin{lstlisting}
def show(self)
\end{lstlisting}
\textbf{Description:} Displays all active Matplotlib plots.

\subsection{\texttt{get\_data}}
\begin{lstlisting}
def get_data(self)
\end{lstlisting}
\textbf{Returns:} The \texttt{self.data} dictionary.

\subsection{\texttt{get\_alpha}}
\begin{lstlisting}
def get_alpha(self)
\end{lstlisting}
\textbf{Returns:} The \texttt{self.\_alpha} dictionary.

\subsection{\texttt{get\_beta}}
\begin{lstlisting}
def get_beta(self)
\end{lstlisting}
\textbf{Returns:} The \texttt{self.\_beta} dictionary.

\subsection{\texttt{operate\_on}}
\begin{lstlisting}
def operate_on(self, data_dict, stat_func, result_attr, *args, **kwargs)
\end{lstlisting}
\textbf{Description:} Applies a statistical function to each column of each table in a data dictionary and stores results.

\textbf{Parameters:}
\begin{itemize}
    \item \texttt{data\_dict}: Dictionary of data tables.
    \item \texttt{stat\_func}: Function to apply to each column.
    \item \texttt{result\_attr}: Attribute name to store results.
    \item \texttt{*args}, \texttt{**kwargs}: Additional arguments for \texttt{stat\_func}.
\end{itemize}

\subsection{\texttt{calc\_distribution}}
\begin{lstlisting}
def calc_distribution(self, bins=50, density=True, *args, **kwargs)
\end{lstlisting}
\textbf{Description:} Calculates histograms for particle size distributions.

\textbf{Parameters:}
\begin{itemize}
    \item \texttt{bins}: Integer or array specifying bin edges (default: 50).
    \item \texttt{density}: Boolean to normalize histogram (default: \texttt{True}).
    \item \texttt{*args}, \texttt{**kwargs}: Additional arguments passed to histogram calculation.
\end{itemize}

\subsection{\texttt{plot\_distribution}}
\begin{lstlisting}
def plot_distribution(self, table_column_map=None)
\end{lstlisting}
\textbf{Description:} Plots particle size distributions with log-normal fits.

\textbf{Parameters:}
\begin{itemize}
    \item \texttt{table\_column\_map}: Optional dictionary mapping table names to column lists (default: \texttt{None}).
\end{itemize}

\textbf{Returns:} List of \texttt{matplotlib.figure.Figure} objects.

\textbf{Notes:} Automatically calculates distributions if not already computed. Supports saving plots if \texttt{self.save\_plot} is \texttt{True}.

\subsection{\texttt{plot\_alphabeta}}
\begin{lstlisting}
def plot_alphabeta(self, *args, **kwargs)
\end{lstlisting}
\textbf{Description:} Generates an $\alpha$ vs. $\beta^2$ scatter plot with regression lines and kinetic regime annotations.

\textbf{Parameters:}
\begin{itemize}
    \item \texttt{*args}: Optional table names or dictionary specifying tables to plot.
    \item \texttt{**kwargs}: Reserved for future extensions.
\end{itemize}

\textbf{Returns:} Tuple of \texttt{matplotlib.figure.Figure} and \texttt{matplotlib.axes.Axes} objects.

\subsection{\texttt{save\_data}}
\begin{lstlisting}
def save_data(self, attr, *args, **kwargs)
\end{lstlisting}
\textbf{Description:} Saves a specified attribute’s data to CSV files.

\textbf{Parameters:}
\begin{itemize}
    \item \texttt{attr}: String name of the attribute to save.
    \item \texttt{*args}, \texttt{**kwargs}: Reserved for future extensions.
\end{itemize}

\subsection{\texttt{export\_data}}
\begin{lstlisting}
def export_data(self, *args, **kwargs)
\end{lstlisting}
\textbf{Description:} Exports specified or all exportable attributes to CSV files.

\textbf{Parameters:}
\begin{itemize}
    \item \texttt{*args}: Optional attribute names to export.
    \item \texttt{**kwargs}: Reserved for future extensions.
\end{itemize}

\section{Usage Examples}
\subsection*{Basic Usage}
\begin{lstlisting}
obj = AlphaBetaSquared('data1.csv', 'data2.csv')
obj.plot_distribution()
obj.plot_alphabeta()
obj.export_data()
\end{lstlisting}

\subsection*{Custom Plotting}
\begin{lstlisting}
obj = AlphaBetaSquared('data.csv')
obj.set_auto_display(False)
figs = obj.plot_distribution({'table1': ['col1', 'col2']})
figs[0].get_axes()[0].set_title('Custom Title')
obj.show()
\end{lstlisting}

\section{Notes}
\begin{itemize}
    \item The \texttt{plot\_config} attribute is partially implemented and marked as a TODO.
    \item Invalid tables or columns in plotting methods are skipped with warnings.
    \item Set \texttt{auto\_display} to \texttt{False} in Jupyter notebooks when manipulating returned figures.
\end{itemize}

\clearpage
\bibliographystyle{apacite}
\bibliography{References}
\listoffigures
\listoftables
\end{document}
